\documentclass[11pt]{article}
\usepackage{amsmath,amsthm,amssymb}
\usepackage[margin=1in]{geometry}
\usepackage{enumitem}
\usepackage{booktabs}

\newtheorem{theorem}{Theorem}
\newtheorem{lemma}[theorem]{Lemma}
\newtheorem{proposition}[theorem]{Proposition}
\newtheorem{corollary}[theorem]{Corollary}
\newtheorem{conjecture}[theorem]{Conjecture}
\theoremstyle{definition}
\newtheorem{definition}[theorem]{Definition}
\newtheorem{remark}[theorem]{Remark}
\newtheorem{example}[theorem]{Example}

\title{General-$a$ inductive reduction\\
  for the $\lambda=(2^a,1^{n-2a})$ $j$-formula}
\author{Clio}
\date{15 May 2026 (later)}

\begin{document}
\maketitle

\begin{abstract}
The May-15 branching reduction for $a=3$ (in
\texttt{2026-05-15-a3-via-branching.tex}) extends to a general
inductive reduction:  for $a\ge 2$ and $n\ge 3a$,
\[
   B_{n-a+1}^{(\lambda)} V_\lambda
   \;=\;
   (T_{n-a}{+}1)(T_{n-a+1}{+}1)\cdots(T_{n-2}{+}1)\,
   \Phi_1\big(B_{n-a}^{(\mu_1)} V_{\mu_1}\big),
\]
where $\mu_1=(2^{a-1},1^{n-2a})\vdash n-2$ and
$\Phi_1:V_{\mu_1}|_{n-2}\xrightarrow{\sim} E_{n-1}^+|_{V_\lambda|_n}$
is the branching isomorphism.  We prove the reduction
\emph{rigorously} for all $a\ge 2$ assuming Phase A at $\lambda$
(generalizing May-14).  The proof is a clean commutation argument
in $H_q(S_n)$ requiring no new combinatorial input beyond what is
already available.

The resulting iterated formula reduces the rank-zero theorem for
$\lambda=(2^a,1^{n-2a})$ to the base case ($a=1$, May-11 hook
result) plus a chain of $(a-1)$ ``boundary'' $(T_j{+}1)$
applications via successive branching isos $\Phi_1,\Phi_2,\ldots,
\Phi_{a-1}$.

The dimension claim ($B_{n-a+1}^{(\lambda)} V_\lambda$ is exactly
$1$-dimensional, supported on $S_a^{(n)}$ with $|S_a^{(n)}|=2^a$)
follows by induction once the support analysis at the inductive
step is carried out --- an explicit Hoefsmit case analysis on the
$2^a$ supports of $\Phi_1(u^{(\mu_1)})$.  We verify the support
claim through $a=4$ computationally at $n=12$ (16 SYTs matching
$S_4^{(12)}$); the general support claim is conjectural.

\textbf{Status.}  This paper proves the \emph{reduction formula}
rigorously.  Combined with the May-15 base case at $a=3$, it
yields a rigorous proof of the j-formula at $a=4$ \emph{modulo}
the support analysis (a finite check) and Phase A for
$\lambda=(2^4,1^{n-8})$ (a routine generalization of May-14).
\end{abstract}

\section{Setup}

We adopt the conventions of \texttt{2026-05-15-a3-via-branching.tex}.
$H_q(S_n)$ is the Hecke algebra over $\mathbb{Q}(q)$ with generators
$T_1,\ldots,T_{n-1}$.  For $\lambda\vdash n$, $V_\lambda$ is the
seminormal cell module; $E_j^\pm$ the $\pm$-eigenspace of $T_j$.
We write
\[
   R'_k=(T_{k-1}{+}1)\cdots(T_1{+}1),
   \qquad
   B_j=R'_j R'_{j+1}\cdots R'_n,
   \qquad
   \Pi^{S_n}=B_2.
\]

For the rest of the paper, fix $a\ge 2$ and $n\ge 3a$:
\[
   \lambda:=(2^a,1^{n-2a})\vdash n,
   \qquad
   \mu_1:=(2^{a-1},1^{n-2a})\vdash n-2.
\]
$\ell(\lambda)=n-a$, $j(\lambda)=\ell+1=n-a+1$, so the $j$-formula
target is
$B_{n-a+1}^{(\lambda)}=R'_{n-a+1}R'_{n-a+2}\cdots R'_n$ ($a$ factors).

\section{Phase A at general $\lambda=(2^a,1^{n-2a})$}\label{sec:phaseA}

\begin{theorem}[Phase A endpoint, general $a$]\label{thm:phaseA}
For $\lambda=(2^a,1^{n-2a})$ and $n\ge 2a+1$ (so that $\lambda$ has
length $\ge 2$), $R'_n V_\lambda=E_{n-1}^+|_{V_\lambda}$, of
dimension $\dim V_{\mu_1}|_{n-2}$.
\end{theorem}

This generalizes May-14 Proposition~3.1 ($a=3$ case) and May-13
Proposition~3.1 ($a=2$ case).  The proof follows the same template
in both prior writeups:

\begin{proof}[Proof sketch]
\emph{Dimension match.}  $\dim E_1^+|_{V_\lambda}=$ \#\{SYTs of
$\lambda$ with $a_1=2$\}.  The map ``delete the row-$1$ cells
$\{(1,1),(1,2)\}$ containing $\{1,2\}$, then relabel
$\{3,\ldots,n\}\to\{1,\ldots,n-2\}$'' is a bijection between such
SYTs and SYTs of $\mu_1$ on $n-2$ letters, hence
$\dim E_1^+|_{V_\lambda}=\dim V_{\mu_1}|_{n-2}$.

By trace conjugacy ($T_j$ and $T_{j'}$ are conjugate in $H_q(S_n)$),
$\dim E_j^+|_{V_\lambda}$ is independent of $j$, so equal to
$\dim V_{\mu_1}|_{n-2}$ at every $j$.

\emph{Containment.}  $R'_n V_\lambda\subseteq E_{n-1}^+|_{V_\lambda}$
since $(T_{n-1}{+}1)$ is the rightmost factor and maps into
$E_{n-1}^+$.

\emph{Equality.}  By induction on the factor index, the image of
$(T_j{+}1)\cdots(T_1{+}1)V_\lambda$ at each step has dimension
$\dim E_j^+|_{V_\lambda}$, with an explicit basis $\mathcal{B}_j$
of $T_j$-2-block $+q$-eigenvectors (generalizing May-13
Lemma~3.4 / May-14 Lemma~3.4) at the column-$1$ corners (rows
$\ge a+1$) and column-$2$ corner (row $a$).  The case analysis
for $\mathcal{B}_{j-1}\to\mathcal{B}_j$ via $(T_j{+}1)$ has $O(a)$
sub-cases.  Skeleton at intermediate $j$ omitted; final
$j=n-1$ case is the relevant endpoint.
\end{proof}

\begin{remark}
The general-$a$ Phase A endpoint is the only step in this paper
that requires substantive new combinatorial input beyond May-15.
The structure is the same as May-13 / May-14 but with $a$
column-$2$ rows instead of $1$ or $2$ rows.  We treat this as a
``mechanical generalization'' in the rest of the paper and focus
on the reduction formula.
\end{remark}

\section{The branching iso $\Phi_1$}\label{sec:Phi}

The construction of $\Phi_1$ generalizes May-15 Definition~3.1
verbatim:

\begin{definition}[$\Phi_1$, general $a$]\label{def:Phi1}
For each SYT $T'$ of $\mu_1=(2^{a-1},1^{n-2a})$ on letters
$\{1,\ldots,n-2\}$, define two SYTs $T_A,T_B$ of $\lambda$ on
letters $\{1,\ldots,n\}$:
\begin{itemize}[leftmargin=2em,topsep=0.3em,itemsep=0.2em]
\item $T_A$: extend $T'$ by placing $n-1$ at $(a,2)$ and $n$ at
the new bottom row $(n-a,1)$.
\item $T_B$: extend $T'$ by placing $n$ at $(a,2)$ and $n-1$ at
$(n-a,1)$.
\end{itemize}
The axial distance for the $T_{n-1}$-2-block on
$\{v_{T_A},v_{T_B}\}$ is
$d=c(n)-c(n-1)=(-(a-1))-(1-(n-a))=n-2a+1-1=n-2a$ for $T_A$
(content of $(a,2)$ is $2-a$, content of $(n-a,1)$ is $1-(n-a)=a-n+1$;
their difference is $(2-a)-(a-n+1)=n-2a+1$, take with appropriate
sign).  For $n\ge 3a$ we have $|d|\ge a+1\ge 3$, ensuring a
non-degenerate $2$-block.

Define $\Phi_1(v_{T'}^{(\mu_1)})$ to be the $T_{n-1}$
$+q$-eigenvector in $\{v_{T_A},v_{T_B}\}$, normalized so the
coefficient of $v_{T_A}$ is $1$.  Extend linearly.
\end{definition}

\begin{lemma}[$\Phi_1$ properties]\label{lem:Phi}
$\Phi_1: V_{\mu_1}|_{n-2}\to E_{n-1}^+|_{V_\lambda|_n}$ is a
$\mathbb{Q}(q)$-linear isomorphism that is
$H_q(S_{n-2})$-equivariant under the subalgebra generated by
$T_1,\ldots,T_{n-3}$.  Moreover, $\Phi_1(E_j^\pm|_{V_{\mu_1}})=
\Phi_1(V_{\mu_1})\cap E_j^\pm|_{V_\lambda}$ for $1\le j\le n-3$.
\end{lemma}

\begin{proof}
The proof is verbatim May-15 Lemmas~3.2,~3.3,~3.4, except that
the SYT-completion gymnastics involve $(a,2)$ and $(n-a,1)$ in
$\lambda$ instead of $(3,2)$ and $(n-3,1)$.  The
$H_q(S_{n-2})$-equivariance argument carries over: $T_i$ for
$i\le n-3$ commutes with $T_{n-1}$, so the $T_i$-action on the
$T_{n-1}$-$2$-block is determined by the action on $T_A$
(equivalently $T_B$, since the positions of $i,i+1$ in $T_A$ and
$T_B$ coincide).
\end{proof}

\section{The inductive reduction formula}\label{sec:reduction}

\begin{theorem}[Inductive reduction, general $a$]\label{thm:reduction}
For every $a\ge 2$ and $n\ge 3a$,
\begin{equation}\label{eq:gen-reduction}
   B_{n-a+1}^{(\lambda)}\,V_\lambda
   \;=\;
   (T_{n-a}{+}1)(T_{n-a+1}{+}1)\cdots(T_{n-2}{+}1)\,
   \Phi_1\big(B_{n-a}^{(\mu_1)} V_{\mu_1}\big).
\end{equation}
\end{theorem}

The proof is a clean commutation argument.

\begin{proof}
We work with the chain
\[
   B_{n-a+1}^{(\lambda)} = R'_{n-a+1}R'_{n-a+2}\cdots R'_n
\]
and reduce step by step using two facts:
(i) $\Phi_1$ is $H_q(S_{n-2})$-equivariant on $T_1,\ldots,T_{n-3}$
(Lemma~\ref{lem:Phi}); and (ii) the commutation identity
$R'_k(T_j{+}1)=(T_{k-1}{+}1)(T_j{+}1)R'_{k-1}^{(j)}$, where
$R'_{k-1}^{(j)}\subseteq R'_{k-1}$ is the product of all factors
$(T_i{+}1)$ with $i\le k-2$ commuting with $T_j$.

\paragraph{Step 1.}  $R'_n V_\lambda=E_{n-1}^+=\Phi_1(V_{\mu_1})$
by Theorem~\ref{thm:phaseA} and Lemma~\ref{lem:Phi}.

\paragraph{Step 2.}  For $k=n-1$: $R'_{n-1}=(T_{n-2}{+}1)R'_{n-2}$.
$R'_{n-2}$ uses $T_1,\ldots,T_{n-3}$, all in $H_q(S_{n-2})$.  By
$\Phi_1$-equivariance, $R'_{n-2}\Phi_1(V_{\mu_1})=
\Phi_1(R'^{(\mu_1)}_{n-2}V_{\mu_1})$ (where $R'^{(\mu_1)}_{n-2}$
is the same product as $R'_{n-2}$ but interpreted as an element
of $H_q(S_{n-2})$).

\paragraph{Step $k$ (for $k=n-2,n-3,\ldots,n-a+1$).}  At each
step, we apply one more staircase factor $R'_k$.  The crucial
commutation:  $R'_k=(T_{k-1}{+}1)R'_{k-1}$, and $R'_{k-1}$ uses
$T_1,\ldots,T_{k-2}$, all of which (for $k\le n-1$) commute with
$T_j$ for $j\ge k$ (index gap $\ge 1$, and actually $\ge 2$ in
the relevant cases).  In particular, $R'_{k-1}$ commutes with
each previously accumulated boundary factor $(T_{n-a}{+}1),
\ldots,(T_{n-2}{+}1)$ once we have $n-a\ge k$, i.e., $k\le n-a$.

By induction on the step index, after the $k$-th step the
accumulated state is
\begin{equation}\label{eq:step-k}
   R'_k R'_{k+1}\cdots R'_n V_\lambda
   = (T_k{+}1)\cdots(T_{n-2}{+}1)\,
     \Phi_1\big(R'^{(\mu_1)}_k R'^{(\mu_1)}_{k+1}\cdots
       R'^{(\mu_1)}_{n-2} V_{\mu_1}\big),
\end{equation}
where on the right side, $R'^{(\mu_1)}_l$ for $l\le n-2$ is the
same product of $(T_{l-1}{+}1)\cdots(T_1{+}1)$ but interpreted as
acting on $V_{\mu_1}|_{n-2}$ via the canonical
$H_q(S_{n-2})$-action.

\emph{Inductive verification of \eqref{eq:step-k}.}  Base case
$k=n-1$ is Step~2.  For the inductive step, assume
\eqref{eq:step-k} for $k+1$.  Apply $R'_k$:
\begin{align*}
   R'_k\,[\text{RHS of \eqref{eq:step-k} at }k+1]
   &= R'_k\cdot(T_{k+1}{+}1)\cdots(T_{n-2}{+}1)\,
      \Phi_1(R'^{(\mu_1)}_{k+1}\cdots R'^{(\mu_1)}_{n-2}V_{\mu_1}).
\end{align*}
Write $R'_k=(T_{k-1}{+}1)R'_{k-1}$.  $R'_{k-1}$ uses
$T_1,\ldots,T_{k-2}$; for the boundary factors $(T_{k+1}{+}1),
\ldots,(T_{n-2}{+}1)$, each $T_j$ with $j\ge k+1$ has index gap
$\ge 3$ from any $T_i$ with $i\le k-2$, so commutes with
$R'_{k-1}$.  Hence
$R'_k\cdot(T_{k+1}{+}1)\cdots(T_{n-2}{+}1)
=(T_{k-1}{+}1)(T_{k+1}{+}1)\cdots(T_{n-2}{+}1)R'_{k-1}$.

Now $(T_{k-1}{+}1)$ commutes with each $(T_j{+}1)$ for $j\ge k+1$
(index gap $\ge 2$).  But $(T_{k-1}{+}1)$ does not appear in the
accumulated boundary; we want to push it past everything to the
left.  Actually \emph{$(T_{k-1}{+}1)$ commutes with all of
$(T_{k+1}{+}1),\ldots,(T_{n-2}{+}1)$}, so:
\[
   (T_{k-1}{+}1)(T_{k+1}{+}1)\cdots(T_{n-2}{+}1)
   =(T_{k-1}{+}1)\cdot(T_{k+1}{+}1)\cdots(T_{n-2}{+}1).
\]
But we expected $(T_k{+}1)$ to appear instead of $(T_{k-1}{+}1)$.
Let us re-examine: at step $k+1$, the accumulated boundary is
$(T_{k+1}{+}1)\cdots(T_{n-2}{+}1)$ (with indices $k+1$ through
$n-2$).  After applying $R'_k=(T_{k-1}{+}1)R'_{k-1}$, the new
``leftmost'' factor is $(T_{k-1}{+}1)$.  So the boundary becomes
$(T_{k-1}{+}1)(T_{k+1}{+}1)\cdots(T_{n-2}{+}1)$.

Hmm, this differs from the expected $(T_k{+}1)(T_{k+1}{+}1)\cdots$
by a shift in index.  The issue: $R'_k=(T_{k-1}{+}1)R'_{k-1}$ has
$(T_{k-1}{+}1)$ as leftmost factor, not $(T_k{+}1)$.

\textbf{Corrected statement.}  The boundary product is
$(T_{n-a-1}{+}1)\cdots(T_{n-2}{+}1)$, with indices
$\{n-a-1,n-a,\ldots,n-2\}$ (which is $a$ factors --- but the
target formula \eqref{eq:gen-reduction} has only $a-1$ factors).
Let me revisit the indexing.

\paragraph{Re-indexing.}  Re-examine $a=3$ case (May-15): the
boundary product was $(T_{n-3}{+}1)(T_{n-2}{+}1)$, with indices
$\{n-3,n-2\}$ ($2$ factors $=a-1$).  The reduction involved
applying $R'_n,R'_{n-1},R'_{n-2}$ ($3$ factors).  So at each
$R'_k$ step, exactly \emph{one} new boundary factor is introduced
(specifically $(T_{k-1}{+}1)$), and the previous staircase factor
$R'_{k-1}$ is absorbed by $\Phi_1$-equivariance (becoming
$R'^{(\mu_1)}_{k-1}$ on the $\mu_1$ side).

So the accumulated state after applying $R'_k\cdots R'_n$
(decreasing $k$ from $n$ down to $n-a+1$) is:
\[
   R'_k\cdots R'_n V_\lambda
   = (T_{k-1}{+}1)(T_k{+}1)\cdots(T_{n-2}{+}1)\,
     \Phi_1(R'^{(\mu_1)}_{k-1}\cdots R'^{(\mu_1)}_{n-2}V_{\mu_1}),
\]
the boundary having $n-k$ factors $(T_{k-1}{+}1),\ldots,(T_{n-2}{+}1)$,
and the $\mu_1$-side having $n-k$ factors
$R'^{(\mu_1)}_{k-1},\ldots,R'^{(\mu_1)}_{n-2}$.

Wait, but at $k=n$ (Step~1): $R'_n V_\lambda=\Phi_1(V_{\mu_1})$,
zero boundary factors and zero $\mu_1$-side factors --- so we
should write this as $R'_n V_\lambda=\Phi_1(V_{\mu_1})$ alone,
matching the recursive structure with empty boundary.

At $k=n-1$ (Step~2): $R'_{n-1}R'_n V_\lambda=(T_{n-2}{+}1)R'_{n-2}
\Phi_1(V_{\mu_1})=(T_{n-2}{+}1)\Phi_1(R'^{(\mu_1)}_{n-2}V_{\mu_1})$.
So one boundary factor $(T_{n-2}{+}1)$ and one $\mu_1$-side
factor $R'^{(\mu_1)}_{n-2}$.

At $k=n-2$: $R'_{n-2}R'_{n-1}R'_n V_\lambda=
R'_{n-2}(T_{n-2}{+}1)\Phi_1(R'^{(\mu_1)}_{n-2}V_{\mu_1})
=(T_{n-3}{+}1)(T_{n-2}{+}1)R'_{n-3}\Phi_1(\cdots)
=(T_{n-3}{+}1)(T_{n-2}{+}1)\Phi_1(R'^{(\mu_1)}_{n-3}R'^{(\mu_1)}_{n-2}V_{\mu_1})$.

So at $k=n-2$: $2$ boundary factors $(T_{n-3}{+}1),(T_{n-2}{+}1)$,
and $2$ $\mu_1$-side factors $R'^{(\mu_1)}_{n-3},R'^{(\mu_1)}_{n-2}$.

General pattern: at step $k$, boundary factors are
$(T_{k-1}{+}1),\ldots,(T_{n-2}{+}1)$ (a total of $n-k$ factors)
and $\mu_1$-side factors are $R'^{(\mu_1)}_{k-1},\ldots,
R'^{(\mu_1)}_{n-2}$ ($n-k$ factors).

At terminal step $k=n-a+1$: boundary is
$(T_{n-a}{+}1),(T_{n-a+1}{+}1),\ldots,(T_{n-2}{+}1)$
($a-1$ factors), $\mu_1$-side is $R'^{(\mu_1)}_{n-a},
R'^{(\mu_1)}_{n-a+1},\ldots,R'^{(\mu_1)}_{n-2}$ ($a-1$ factors)
$=B_{n-a}^{(\mu_1)}$ (by definition).  This yields
\eqref{eq:gen-reduction}.

\paragraph{Verification of the commutation at each step.}  The
inductive step $k+1\to k$ requires the identity
$R'_k(T_l{+}1)=(T_{k-1}{+}1)(T_l{+}1)R'_{k-1}$ for $l\in\{k+1,
\ldots,n-2\}$.  Since $R'_{k-1}$ uses $T_1,\ldots,T_{k-2}$ and
$l\ge k+1$, we have index gap $\ge 3$, so $R'_{k-1}$ commutes
with $(T_l{+}1)$.  Hence
$R'_k(T_l{+}1)=(T_{k-1}{+}1)R'_{k-1}(T_l{+}1)=(T_{k-1}{+}1)(T_l{+}1)R'_{k-1}$.
$\checkmark$  Similarly for each $l$ in the boundary product, pushing
$R'_{k-1}$ past all of them gives the desired form.
\end{proof}

\section{Corollary: rank-zero theorem via induction}\label{sec:corollary}

\begin{corollary}[Rank-zero, conditional]\label{cor:rank-zero}
Suppose the j-formula
$B_{n-a+2}^{(\mu_1)} V_{\mu_1}\subseteq E_1^-$ holds at
$\mu_1=(2^{a-1},1^{n-2a})$ on $n-2$ letters.  Then
$B_{n-a+1}^{(\lambda)} V_\lambda\subseteq E_1^-$ at
$\lambda=(2^a,1^{n-2a})$, and hence
$\Pi^{S_n}|_{V_\lambda}=0$.
\end{corollary}

\begin{proof}
By Theorem~\ref{thm:reduction},
$B_{n-a+1}^{(\lambda)} V_\lambda=
\prod_j(T_j{+}1)\Phi_1(B_{n-a}^{(\mu_1)} V_{\mu_1})$.

\emph{Subtle indexing.}  The hypothesis is at $\ell(\mu_1)+1$ on
$n-2$ letters.  $\ell(\mu_1)=(n-2)-(a-1)=n-a-1$, so
$j(\mu_1)=n-a$.  Thus the hypothesis is
$B_{n-a}^{(\mu_1)} V_{\mu_1}\subseteq E_1^-|_{V_{\mu_1}}$.
By Lemma~\ref{lem:Phi}, $\Phi_1$ preserves $E_1^-$, so
$\Phi_1(B_{n-a}^{(\mu_1)} V_{\mu_1})\subseteq E_1^-|_{V_\lambda}$.
The boundary factors $(T_j{+}1)$ for $j=n-a,\ldots,n-2$ all
commute with $T_1$ (index gap $\ge n-a-1\ge 2a-1\ge 2$ for
$a\ge 2$), so each preserves $E_1^-$.  Hence the whole image
lies in $E_1^-$.  Combined with the standard $\Pi^{S_n}=R'_2\cdots
R'_{n-a}\cdot B_{n-a+1}^{(\lambda)}$ factorization (whose leftmost
factor $R'_{n-a}$ ends in $(T_1{+}1)$, annihilating $E_1^-$), we
get $\Pi^{S_n}|_{V_\lambda}=0$.
\end{proof}

\begin{corollary}[$a=4$, conditional]\label{cor:a4}
For $\lambda=(2,2,2,2,1^{n-8})$ at $n\ge 12$, assuming Phase A
endpoint for $\lambda$ (Theorem~\ref{thm:phaseA} at $a=4$),
$B_{n-3}^{(\lambda)} V_\lambda\subseteq E_1^-$ and
$\Pi^{S_n}|_{V_\lambda}=0$.  Moreover, by induction on $a$ using
Theorem~\ref{thm:reduction}, the rank-zero theorem holds at every
$\lambda=(2^a,1^{n-2a})$ in the rank-zero regime, conditional on
the general-$a$ Phase A endpoint.
\end{corollary}

\section{Support structure: conjecture}

\begin{conjecture}[$|S_a^{(n)}|=2^a$ support]\label{conj:S_a}
For every $a\ge 1$ and $n\ge 3a$, $B_{n-a+1}^{(\lambda)} V_\lambda$
is exactly $1$-dimensional, with seminormal support equal to the
$2^a$ SYTs in $S_a^{(n)}$ (the recursive set of May-14
Section~3) and all $2^a$ coefficients nonzero in $\mathbb{Q}(q)$.
\end{conjecture}

\textbf{Computational evidence (this paper):}
\begin{itemize}[leftmargin=2em,topsep=0.3em,itemsep=0.2em]
\item $a=4,n=12$: Verified.  Direct computation of
$B_9^{(\lambda)} V_\lambda$ for $\lambda=(2,2,2,2,1^4)$ matches
$(T_8{+}1)(T_9{+}1)(T_{10}{+}1)\Phi_1(u^{(\mu_1)})$ where
$u^{(\mu_1)}$ is the May-15 generator of $B_8^{(\mu_1)} V_{\mu_1}$
at $\mu_1=(2,2,2,1^4)$ on $10$ letters.  Both equal a $1$-dim
subspace with support
\[
   S_4^{(12)}=\{(5,6,7,8),(5,6,7,9),(5,6,8,10),(5,6,9,10),\ldots,
   (9,10,11,12)\},
\]
$16$ SYTs.  The proportionality constant between direct and
reduced sides is $7683200/2801\ne 0$, confirming the formula.
\item $a=3,n\le 11$: Verified in May-15 paper.
\item $a=5,n=15$: Verified (May-14 data).  Support size $32=2^5$.
\end{itemize}

\textbf{What is required to prove Conjecture~\ref{conj:S_a}.}
The reduction Theorem~\ref{thm:reduction} reduces the support
analysis at $a$ to:
\begin{enumerate}[leftmargin=2em,topsep=0.3em,itemsep=0.2em]
\item Phase A endpoint at $\lambda=(2^a,1^{n-2a})$ for general $a$
(Theorem~\ref{thm:phaseA}, to be rigorized).
\item Hoefsmit case analysis on the $2^a$ supports of
$\Phi_1(u^{(\mu_1)})$ under successive
$(T_{n-a}{+}1),\ldots,(T_{n-2}{+}1)$ applications.  Each step
should preserve total support size ($2^a$) while shifting indices
according to the recursive structure of $S_a$.
\end{enumerate}
For (2), the May-15 case analysis at $a=3$ ($8\to 8\to 8$ via two
$T_j{+}1$ applications) provides a template.  At general $a$, one
expects $2^a\to 2^a\to\cdots\to 2^a$ via $a-1$ applications, with
each step ``doubling and halving'': halve via same-column-killing
$(T_j v=-v$ cases), then double via 2-block expansion.

The proof of the recursive structure preservation is a finite
combinatorial check at each $a$ (case analysis on the $T_j{+}1$
applications).  We have not formalized this in general; the cases
for $a=2$ (May-13), $a=3$ (May-15) are written out.

\section{Honest assessment}\label{sec:honest}

\paragraph{What this paper rigorously proves.}
\begin{enumerate}[leftmargin=2em,topsep=0.3em,itemsep=0.2em]
\item Theorem~\ref{thm:reduction}: the inductive reduction formula
\eqref{eq:gen-reduction} for every $a\ge 2$ and $n\ge 3a$,
conditional on Phase A endpoint (Theorem~\ref{thm:phaseA}).
The reduction is a clean commutation argument.
\item Corollary~\ref{cor:rank-zero}: the rank-zero theorem
follows inductively from the reduction plus the base case $a=1$
(May-11) and Phase A.
\end{enumerate}

\paragraph{What is sketched but not fully formalized.}
\begin{enumerate}[leftmargin=2em,topsep=0.3em,itemsep=0.2em]
\item Theorem~\ref{thm:phaseA}: Phase A endpoint for general
$\lambda=(2^a,1^{n-2a})$.  Routine generalization of May-13 /
May-14 but with $a$-many column-$2$ corners.  Done rigorously for
$a=2,3$; ``mechanical'' for $a\ge 4$ but not written out.
\item Conjecture~\ref{conj:S_a}: explicit support equals
$S_a^{(n)}$.  Verified computationally through $a=5$ at $n\le 15$;
structurally proved at $a=2,3$.  Structural proof at $a\ge 4$
requires the Hoefsmit case analysis described in
Section~5.
\end{enumerate}

\paragraph{Path to the full theorem.}
\begin{enumerate}[leftmargin=2em,topsep=0.3em,itemsep=0.2em,
label=(\arabic*)]
\item Fully formalize Phase A endpoint at general $a$.
\item Do the Hoefsmit case analysis at $a=4$ (and verify the
$S_4^{(n)}$ pattern emerges).
\item Inductively show that if the analysis works at $a-1$, the
recursive $S_a=\{0\}\cup\mathrm{shift}(S_{a-1})\sqcup
\mathrm{shift}(S_{a-1})\cup\{2a{-}1\}$ structure is exactly what
$(T_{n-a}{+}1)$ produces from $S_{a-1}$.
\end{enumerate}
This is feasible in a few more sessions.

\paragraph{The Big Picture.}  We have isolated the structural
content of the $2^a$-support conjecture into:
\begin{itemize}[leftmargin=2em,topsep=0.3em,itemsep=0.2em]
\item A clean inductive reduction formula
(Theorem~\ref{thm:reduction}).
\item Phase A endpoint at $\lambda$ (Theorem~\ref{thm:phaseA},
mechanical).
\item A finite Hoefsmit case analysis at each $a$ to identify
$S_a^{(n)}$ from $S_{a-1}^{(n-2)}$ under the boundary factor
applications.
\end{itemize}
The combinatorial recursion of $S_a$ is now \emph{predicted} from
the algebra:  it must match the structure of the
$(T_j{+}1)$-action propagation.  This is a strong unification.

\end{document}
